% SPDX-FileCopyrightText: Copyright (c) 2016-2026 Objectionary.com
% SPDX-License-Identifier: MIT

\section{Outer Kinds Reference}\label{app:outer-kinds}

This appendix is an at-a-glance reference for the twelve outer kinds of
  EO expressions.
The conceptual treatment --- what an outer kind is, why it matters, and
  how the openness state machine works --- is in
  \cref{sec:syntax:outer-kinds}.
The per-kind table below collects two structural properties that govern
  how each construct may be extended in the surrounding line stream:
  whether the construct admits \emph{deeper-indent children},
  and whether it may be \emph{wrapped} by a same-indent
  \ff{.method} continuation.

Below, \enquote{after block closes} means the kind's openness has
  transitioned from \emph{open} to \emph{vertical-completed} ---
  i.e., its deeper-indent body block has ended but a \ff{.method}
  wrap is still permissible (\cref{sec:wf:notation}).

The column \emph{Cardinality} records how the construct occupies the
  source:
  \enquote{1 line} means the entire expression lives on a single source
  line with no deeper-indent block;
  \enquote{1 line + body} means a head line followed by a deeper-indent
  body block;
  \enquote{multi-line} means a multi-line construct without a single
  syntactic head (a method chain, a multi-line text block, or a vertical
  application).
The column \emph{Children} answers \enquote{may receive deeper-indent
  children?} and \emph{Wraps} answers \enquote{may be wrapped by a
  same-indent \ff{.method} continuation?}.
A bare \enquote{yes} means the construct is always open in that
  direction; \enquote{yes (after)} means only after the body block
  closes; \enquote{no} means never.

\begin{table*}
\capt{The twelve outer kinds of EO expressions.}
\label{tab:outer-kinds}
\footnotesize
\begin{tabularx}{\textwidth}{l l c c X}
\toprule
Outer kind & Cardinality & Children & Wraps & Notes \\
\midrule
\ff{head}                 & 1 line        & yes        & yes (after) &
  Bare identifier, literal, paren group, or \ff{*}; no chain, no hargs. \\
\ff{hmethod}              & 1 line        & yes\(^*\)  & yes (after) &
  \ff{x.y.z} chained dispatch on one line. \(^*\)Children admitted only when the chain has zero horizontal args. \\
\ff{happlication}         & 1 line        & no         & no &
  Head + \(\geq\)1 horizontal args. Subsumes a chained head; the chain is inside, the outer kind is happlication. \\
\ff{bare-formation}       & 1 line + body & yes        & yes (after) &
  \ff{[params] [> name]}. Body is a deeper-indent block of children. \\
\ff{bare-reversed}        & 1 line + body & yes        & yes (after) &
  \ff{name.} vertical form. First deeper child is the receiver; siblings are method args. \\
\ff{reversed-with-hargs}  & 1 line        & no         & no &
  Horizontal form \ff{name. recv args\ldots}. \\
\ff{compact-tuple}        & 1 line + body & yes        & yes (after) &
  \ff{head *N}. First \(N\) children are direct positional args; the rest tupled into a \(\Phi.\mathtt{tuple}\) wrapper. \\
\ff{inline-phi-formation} & 1 line        & no         & no &
  \ff{expr > [params] > name}. LHS becomes the \(\varphi\)-body of the inline formation. \\
\ff{vapplication}         & multi-line    & yes        & yes (after) &
  Head line plus a deeper-indent block of vertical arguments. \\
\ff{vmethod}              & multi-line    & yes\(^*\)  & yes\(^\dagger\) &
  Chain of same-indent \ff{.method} continuations. \(^*\)Children allowed only if the last \ff{.method} has 0 hargs; \(^\dagger\)more \ff{.method}s extend the chain. \\
\ff{vmethod-with-hargs}   & multi-line    & no         & no &
  A \ff{.method} continuation carrying \(\geq\)1 horizontal args; closes the chain immediately. \\
\ff{text-block}           & multi-line    & ---        & yes (after) &
  \ff{"""}-delimited body. Body lines are lexical content, not children. \\
\bottomrule
\end{tabularx}
\end{table*}

Four of the twelve outer kinds are \emph{horizontally-completed}:
  they admit no deeper children and cannot be wrapped by a same-indent
  \ff{.method} continuation.
The set is

\begin{equation*}
  \mathcal{H} = \{\,
    \mathtt{happlication},\
    \mathtt{reversed\text{-}with\text{-}hargs},\
    \mathtt{vmethod\text{-}with\text{-}hargs},\
    \mathtt{inline\text{-}phi\text{-}formation}
    \,\}.
\end{equation*}

The well-formedness rules of \cref{sec:wf:notation} are stated in terms of
  membership in \(\mathcal{H}\).
